\paragraph{临界预算的熵解释与 Zeckendorf 基准分解}
本段把临界斜率 \(a_1=\tau'(1^-)\) 与临界预算 \(\beta_{\mathrm{crit}}=a_1/\log 2\) 解释为“折叠推前分布的条件熵亏损率”，
并将其精确拆分为 Zeckendorf 结构基准项与非均匀性惩罚项。

\begin{theorem}[临界预算的熵恒等式]\label{thm:pom-oracle-critical-budget-entropy-identity}
\leanverified{paper\_pom\_oracle\_critical\_budget\_entropy\_identity}
令 \(w_m(x)=d_m(x)/2^m\) 为均匀微态推前分布，并定义其 Shannon 熵
\[
H_1(m):=-\sum_{x\in X_m}w_m(x)\log w_m(x).
\]
若 \(\tau\) 在 \(q=1\) 的左邻域可微并满足 \(a_1:=\tau'(1^-)\)（同定理 \ref{thm:pom-oracle-critical-window-gaussian}），
则熵率极限
\[
h_1:=\lim_{m\to\infty}\frac{1}{m}H_1(m)
\]
存在且满足严格恒等式
\[
\boxed{\;
h_1=\log 2-a_1,\qquad
\beta_{\mathrm{crit}}=1-\frac{h_1}{\log 2}\; } .
\]
\end{theorem}
\begin{proof}
由 \(w_m(x)=d_m(x)/2^m\) 得
\[
H_1(m)
=-\sum_x \frac{d_m(x)}{2^m}\log\frac{d_m(x)}{2^m}
=m\log 2-\sum_x w_m(x)\log d_m(x).
\]
另一方面，
\[
\frac{\dd}{\dd q}\log S_q(m)\Big|_{q=1}
=\frac{\sum_x d_m(x)\log d_m(x)}{\sum_x d_m(x)}
=\sum_x w_m(x)\log d_m(x).
\]
在 \(\tau\) 的可微性制度下，\(\frac{1}{m}\log S_q(m)\to\tau(q)\) 且在 \(q=1^-\) 可交换求导，
从而
\[
\lim_{m\to\infty}\frac{1}{m}\sum_x w_m(x)\log d_m(x)
=\tau'(1^-)=a_1.
\]
将其代回 \(H_1(m)=m\log 2-\sum_x w_m(x)\log d_m(x)\) 并除以 \(m\) 即得 \(h_1=\log 2-a_1\)；
再代入 \(\beta_{\mathrm{crit}}=a_1/\log 2\) 得第二式。证毕。
\end{proof}

\begin{corollary}[极端固定而熵可调：在非冻结相中 \texorpdfstring{$(\alpha_*,g_*)$}{(alpha*,g*)} 不决定 \texorpdfstring{$\beta_{\mathrm{crit}}$}{beta\_crit}]\label{cor:pom-extremes-do-not-determine-beta-crit}
设 \(d_m:X_m\to\NN_{\ge 1}\) 为一族纤维谱，并记
\[
D_m:=\max_{x\in X_m}d_m(x),\qquad
K_m:=\card{\{x\in X_m:\ d_m(x)=D_m\}}.
\]
假设极限
\[
\alpha_*:=\lim_{m\to\infty}\frac{1}{m}\log D_m,\qquad
g_*:=\lim_{m\to\infty}\frac{1}{m}\log K_m
\]
存在且满足 \(\alpha_*+g_*<\log2\)。
并假设类型空间具有指数规模：存在常数 \(c_X>0\) 使 \(\log|X_m|\ge c_X m\) 对充分大 \(m\) 成立，且 \(c_X>\log2-\alpha_*\)。
则在保持 \((D_m,K_m)\)（从而保持 \((\alpha_*,g_*)\)）不变的前提下，仅通过重新分配非最大层质量即可构造两族谱 \(\{d_m\}\)、\(\{\tilde d_m\}\)，使得对应推前分布的熵率 \(h_1\) 收敛到两个不同常数；从而由定理 \ref{thm:pom-oracle-critical-budget-entropy-identity}，\(\beta_{\mathrm{crit}}=1-h_1/\log2\) 亦可在一段正长度区间内变化。
\end{corollary}
\begin{proof}
由 \(\alpha_*+g_*<\log2\) 得
\[
\frac{K_mD_m}{2^m}=\exp\!\bigl((\alpha_*+g_*-\log2+o(1))m\bigr)\ \longrightarrow\ 0,
\]
故最大层在概率意义下的质量占比趋于零，而剩余质量
\(
R_m:=2^m-K_mD_m
\)
满足 \(R_m=2^m(1-o(1))\)。
\par
取常数 \(c_-,c_+\) 使 \(\log2-\alpha_*<c_-<c_+<c_X\)。
对每个 \(m\)，令 \(L_m^\pm:=\lfloor e^{c_\pm m}\rfloor\) 并取互异类型点集 \(Y_m^\pm\subseteq X_m\) 满足 \(|Y_m^\pm|=L_m^\pm\) 且与最大层达到者集合不交（这由 \(c_\pm<c_X\) 保证可行）。
定义两族谱：在两种构造中都令最大层上恰有 \(K_m\) 个类型取值 \(D_m\)，并将剩余质量 \(R_m\) 尽量均匀分配到 \(Y_m^\pm\) 上（再作整数修正与余量收尾），其余类型取值为 \(1\)。
由于
\[
\frac{1}{m}\log\Bigl(\frac{R_m}{L_m^\pm}\Bigr)
=(\log2-c_\pm)+o(1)\ <\ \alpha_*,
\]
对充分大 \(m\) 可保证分配到 \(Y_m^\pm\) 的每个数值严格小于 \(D_m\)，从而两族谱具有相同的 \((D_m,K_m)\)。
\par
令 \(w_m,\tilde w_m\) 为两族谱对应的推前分布。
在两种构造下，除去最大层的 \(o(1)\) 质量后，剩余质量以 \(1-o(1)\) 的比例近似均匀分布在 \(L_m^\pm\) 个类型上，
因此 Shannon 熵满足
\[
H_1(m)=\log L_m^\pm+o(m),
\]
从而熵率极限分别为
\(
h_1=\lim\frac1m H_1(m)=c_\pm
\)，且 \(c_-\ne c_+\)。
最后由定理 \ref{thm:pom-oracle-critical-budget-entropy-identity} 的恒等式 \(\beta_{\mathrm{crit}}=1-h_1/\log2\) 得所示结论。证毕。
\end{proof}

\begin{theorem}[Zeckendorf 基准项与熵缺口率分解]\label{thm:pom-oracle-critical-budget-zeckendorf-plus-gap}
\leanverified{paper\_pom\_oracle\_critical\_budget\_zeckendorf\_plus\_gap}
令 \(U_m=\mathrm{Unif}(X_m)\) 为类型均匀分布，并设极限
\[
\gamma:=\lim_{m\to\infty}\frac{1}{m}D_{\mathrm{KL}}(w_m\|U_m)
=\lim_{m\to\infty}\frac{1}{m}\bigl(\log|X_m|-H_1(m)\bigr)
\]
存在。
若进一步有 \(\lim_{m\to\infty}\frac{1}{m}\log|X_m|=\log\varphi\)，则
\[
\boxed{\;
a_1=(\log 2-\log\varphi)+\gamma,\qquad
\beta_{\mathrm{crit}}=1-\frac{\log\varphi}{\log 2}+\frac{\gamma}{\log 2}\; } .
\]
\end{theorem}
\begin{proof}
由恒等式 \(D_{\mathrm{KL}}(w_m\|U_m)=\log|X_m|-H_1(m)\) 得
\(
\gamma=\log\varphi-h_1
\)。
将定理 \ref{thm:pom-oracle-critical-budget-entropy-identity} 的 \(a_1=\log 2-h_1\) 与上式合并，
得到 \(a_1=(\log 2-\log\varphi)+\gamma\)，再除以 \(\log 2\) 即得对 \(\beta_{\mathrm{crit}}\) 的分解。证毕。
\end{proof}

\endinput
